% Generated draft: requires mathematical and source review.
\documentclass[11pt]{article}
\usepackage[margin=25mm]{geometry}
\usepackage{fontspec}
\setmainfont{Latin Modern Roman}
\usepackage{amsmath,amssymb,mathtools,unicode-math}
\setmathfont{Latin Modern Math}
\usepackage{enumitem,needspace,xurl,hyperref}
\setlength{\parindent}{0pt}
\setlength{\parskip}{5pt}
\setlength{\emergencystretch}{3em}
\newcommand{\N}{\mathbb{N}}
\newcommand{\Z}{\mathbb{Z}}
\newcommand{\Q}{\mathbb{Q}}
\newcommand{\R}{\mathbb{R}}
\newcommand{\C}{\mathbb{C}}
\newcommand{\F}{\mathbb{F}}
\newcommand{\A}{\mathbb{A}}
\newcommand{\PP}{\mathbb{P}}
\newcommand{\E}{\mathbb{E}}
\newcommand{\eps}{\varepsilon}
\newcommand{\dd}{\,\mathrm{d}}
\begin{document}
\section*{UnsolvedMath ID: 701513}
\textbf{Source problem code:} catalogue entry
\subsection{Definitions and mathematical statement}
Let \(\pi\) be any set of primes. A finite group \(H\) is a \(\pi\)-group if every prime divisor of \(|H|\) belongs to \(\pi\); the trivial group is included. A \(\pi\)-subgroup of a finite group is a subgroup that is a \(\pi\)-group. For a subset \(X\) of a group, write \(\langle X\rangle\) for the subgroup generated by \(X\).

A finite group \(G\) is a \(BS_{\pi}\)-group if, for every conjugacy class \(C\) of \(G\),
\[
\left(\forall x,y\in C,\quad \langle x,y\rangle\text{ is a }\pi\text{-group}\right)
\quad\Longrightarrow\quad
\langle C\rangle\text{ is a }\pi\text{-group}.
\]
\textbf{Question (D. O. Revin).} For every set of primes \(\pi\), every finite \(BS_{\pi}\)-group \(G\), and every normal subgroup \(N\trianglelefteq G\), must \(N\) be a \(BS_{\pi}\)-group? Explicitly, must every conjugacy class \(D\) of \(N\) satisfy
\[
\left(\forall x,y\in D,\quad \langle x,y\rangle\text{ is a }\pi\text{-group}\right)
\quad\Longrightarrow\quad
\langle D\rangle\text{ is a }\pi\text{-group}?
\]
Here conjugacy classes of \(N\) are taken under conjugation by elements of \(N\).

The supplied source reports that the answer is affirmative when \(2\notin\pi\), citing D. O. Revin, \emph{Siberian Math. J.} \textbf{52}, no.~2 (2011), 340--347. This qualification is retained as a supplied literature assertion and has not been independently verified.
\subsection{Short English statement}
A finite \(BS_{\pi}\)-group has the property that any conjugacy class whose pairs generate \(\pi\)-subgroups also generates a \(\pi\)-subgroup as a whole. Revin asks whether this property passes to every normal subgroup, using that subgroup's own conjugacy classes. The supplied record labels the problem open and reports an affirmative answer when \(2\notin\pi\).
\subsection{Sources}
\begin{itemize}
\item Local UnsolvedMath registry, permanent ID \(701513\), exact source code \texttt{catalogue entry}, title \emph{Kourovka Notebook Problem 18.84}. The supplied record labels the problem \emph{open}; this status has not been independently verified.
\item \emph{The Kourovka Notebook (Unsolved Problems in Group Theory)}, problem 18.84, proposed by D. O. Revin in 2014. The supplied background identifies the 18th issue (2014) and the 21st issue edition (2026), arXiv:1401.0300v48. The supplied collection description names E. I. Khukhro and V. D. Mazurov as editors.
\item D. O. Revin, \emph{Siberian Math. J.} \textbf{52}, no.~2 (2011), 340--347. The supplied statement cites this reference for the affirmative answer when \(2\notin\pi\); no article title is supplied, and the reference and its asserted implication have not been independently checked.
\end{itemize}
\end{document}
